\chapter[Menschenverherrlichung]{Menschenverherrlichung}
\quellseite{8}
Im October war der Allgemeine Gustav-Adolf-Verein in Hamburg versammelt, der sich ganz besonders die Vertheidigung des evangelischen Bekenntnisses gegen das Papstthum zum Zweck gesetzt hat. Am Schluß der Versammlung machten die Theilnehmer eine Fahrt zum Grabe Bismarcks. Sie bildeten einen Halbkreis vor dem Eingang, und dann trat der Oberconsistorialrath Superintendent Dr. Dibelius aus Dresden, ein Vertreter der lutherischen Landeskirche Sachsens, vor die halbgeöffnete Grabesthür und sprach, nach innen gerichtet:

\begin{verse}
Hier ruht in Frieden, gewaltiger Held,\\
Deines Geistes irdisches Wandergelt.\\
Du aber stirbst uns nimmermehr,\\
Du, deutschen Volkes Ehr und Wehr.\\
Dein Geist entflammt uns stets aufs neu\\
Zu Gottesfurcht und Mannestreu,\\
Zu wahrer als unser heiliges Gut,\\
Was du erstrittest mit Eisen und Blut.\\
Wir wollen deutsch sein bis ins Mark;\\
Das Evangelium mach uns stark!\\
Das Vaterland durchglüht sein:\\
Du schaust vom Himmel segnend drein!
\end{verse}

Das ist wirklich stark. Mit Recht wird daher diesen vorgeblichen „Bekämpfern“ des Papstthums von der „Hannoverschen Pastoral-Correspondenz“ vorgehalten, daß sie nichts anderes gethan als einen Mißbrauch der römischen Kirche nachgeahmt und Bismarck als einen protestantischen „Heiligen“ verehrt haben. Der ganze Vers ist un-evangelisch und unchristlich. — Vor Kurzem starb in Berlin der berühmte Geschichtsforscher Prof. Theodor Mommsen und wurde dann mit großem Prunke und kirchlichen Ehrenbezeigungen begraben. Die Leiche wurde in die Kaiser-Wilhelm-Gedächtnißkirche gebracht und dort vor dem mit Lichtern umgebenen Altar aufgefahret, was sonst in Berlin nicht geschieht. Daß Mommsens Stellung zum Christenthum eine sehr liberale war, läßt sich schon daraus erkennen, daß er in seinem Testament bestimmte, daß kein anderer als der bekannte irreligiöse Prof. Harnack an seinem Sarge das Wort ergreifen dürfe. Harnack bestieg auch die Kanzel und sprach über das Schriftwort: „Ich habe euch gesetzt, daß ihr hingehet und Frucht bringet, und eure Frucht bleibe.“ Der „Freimund“ sagt: „Die Frucht, die dem Verstorbenen nachgerühmt wurde, waren seine wissenschaftlichen Leistungen. Diese sind natürlich bald überholt und vergessen. Der Herr Christus meint eine ganz andere Frucht, die für die Ewigkeit gewirkt wird. Vom geistlichen Leben und vom Leben der zukünftigen Welt ist in der ganzen Ansprache mit keinem Wort die Rede. ‚Gott‘ wurde nur einmal im Schlußsatz genannt: ‚Wir aber bitten Gott, daß er unsere Arbeit segnen möge wie die Theodor Mommsens.‘“ — Es ist aber ein Mißbrauch des Gotteshauses, wenn es so zur Menschenvergötterung gebraucht wird. L. F.
